\chapter{Component Architecture and Design}

\section{Layered Design}

The gateway follows a compact layered design. Each layer has one dominant
responsibility so that request handling, business rules, data access, and
persistent state remain understandable as separate concerns.

The stack also groups naturally into three bands. The upper band accepts
requests and enforces access boundaries, the middle band converts validated
requests into business decisions, and the lower band anchors those decisions in
durable entities and long-lived state. Reading the architecture this way makes
it clear where interface concerns stop and where business consistency begins.

The visual stack in Figure 5.1 appears after the responsibility map so the
reader can connect the abstract layers to their concrete roles in the system.

\section{Component Responsibilities}

\begin{longtable}{L{0.26\textwidth}L{0.64\textwidth}}
\caption{Component responsibilities}\\
\toprule
\textbf{Component} & \textbf{Responsibility} \\
\midrule
\endfirsthead
\toprule
\textbf{Component} & \textbf{Responsibility} \\
\midrule
\endhead
Internal operations dashboard & Supports merchant review, credential management, account administration, transaction investigation, reconciliation, audit inspection, and recovery actions. \\
\bodyrowrule
Merchant dashboard & Supports merchant-scoped overview, analytics, transaction exploration, credential visibility, and account self-service. \\
\bodyrowrule
Interface layer & Accepts incoming requests, applies access boundaries, validates high-level request intent, and shapes outward responses. \\
\bodyrowrule
Request and response contracts & Define the information each interaction accepts or returns so that system behavior remains predictable. \\
\bodyrowrule
Application services & Enforce merchant readiness, payment and refund rules, notification behavior, analytics preparation, reconciliation, and audit-aware state changes. \\
\bodyrowrule
Data access layer & Retrieves domain data, writes business state changes, and prepares filtered or aggregated views for dashboards and review workflows. \\
\bodyrowrule
Domain entities and rules & Define states, relationships, constraints, and invariants for merchants, transactions, notifications, accounts, and review records. \\
\bodyrowrule
Persistent storage & Holds durable operational and business evidence across all major workflows. \\
\bottomrule
\end{longtable}

\begin{figure}[!htbp]
\centering
\scalebox{0.72}{%
\begin{tikzpicture}[node distance=8mm]
  \node[box, text width=44mm] (interface) {Interface layer\\access control and request handling};
  \node[box, below=of interface, text width=44mm] (contracts) {Request and response contracts\\validated input and output structures};
  \node[box, below=of contracts, text width=44mm] (services) {Application services\\business rules and workflow control};
  \node[box, below=of services, text width=44mm] (dataaccess) {Data access layer\\lookup, list, and aggregation operations};
  \node[box, below=of dataaccess, text width=44mm] (domain) {Domain entities and rules\\states, relationships, and invariants};
  \node[db, below=of domain, text width=36mm] (storage) {Persistent\\storage};

  \draw[arrow] (interface) -- (contracts);
  \draw[arrow] (contracts) -- (services);
  \draw[arrow] (services) -- (dataaccess);
  \draw[arrow] (dataaccess) -- (domain);
  \draw[arrow] (domain) -- (storage);
\end{tikzpicture}%
}
\caption{Component layer stack. Higher layers coordinate interaction, while
lower layers preserve business state and data consistency.}
\end{figure}

\section{Business Rules And State Models}

The domain behavior is organized around explicit state models instead of
implicit status changes hidden inside request handlers.

\begin{table}[!htbp]
\centering
\caption{Core state models}
\begin{tabularx}{\textwidth}{L{0.24\textwidth}L{0.24\textwidth}Y}
\toprule
\textbf{Domain object} & \textbf{Persisted states} & \textbf{Design note} \\
\midrule
Merchant operational status & Pending review, active, rejected, suspended, disabled & Merchant readiness and merchant operating status are related but not identical, so review outcomes and live operating state are tracked separately. \\
\bodyrowrule
Payment transaction & Pending, successful, failed, expired & Payment requests remain pending until confirmed by external evidence or time-based expiration handling. \\
\bodyrowrule
Refund transaction & Refund pending, refunded, refund failed & The system implements full refunds only, and final outcomes depend on external refund evidence. \\
\bodyrowrule
Webhook event & Pending, delivered, failed & Merchant notification delivery is tracked independently from payment and refund finality. \\
\bottomrule
\end{tabularx}
\end{table}

Two design choices are especially important:

\begin{itemize}
  \item payment expiration creates a durable business outcome instead of simply
  removing an old pending request;
  \item ambiguous external evidence produces a reconciliation record instead of
  forcing a risky automatic correction.
\end{itemize}

\section{Asynchronous Event Lifecycle And Recovery}

Although the project does not yet rely on a permanently running background
worker, it already contains the core asynchronous control functions needed for
time-based and retry-based behavior.

\begin{enumerate}
  \item A merchant payment request creates a pending payment with expiration
  metadata.
  \item External payment or refund evidence is stored before the system applies
  state transitions.
  \item Final money-flow outcomes create merchant notification events.
  \item Scheduled processing expires overdue pending payments and continues
  delivery attempts for due notifications.
  \item Delivery failures and ambiguous outcomes are preserved as operational
  evidence rather than being hidden.
  \item Internal users can retry eligible failed notifications without
  changing already-final payment or refund outcomes.
\end{enumerate}

This separation keeps the architecture testable. Core business decisions remain
inside focused application services, while any future scheduling layer can stay
thin and orchestration-oriented.
